\documentclass[11pt]{article}
\usepackage[margin=1in]{geometry}
\usepackage{amsmath,amssymb}
\usepackage[hidelinks]{hyperref}
\usepackage{microtype}

\title{The Root-System Extension of Ramanujan's Master Theorem\\
Was Answered by arXiv:1211.0024}
\author{Literature-status note for arXiv:1103.5126\\
Run \texttt{fa\_banach\_001}, lane 6}
\date{11 August 2026}

\begin{document}
\maketitle

\begin{center}
\fbox{\parbox{0.9\textwidth}{\textbf{Status: literature already answered.}
The later paper is by the same authors, explicitly identifies the earlier
symmetric-space theorem, and proves the requested root-system extension for
arbitrary positive multiplicities.}}
\end{center}

\section*{Source question}

Gestur Ólafsson and Angela Pasquale,
\emph{Ramanujan's Master Theorem for Riemannian symmetric spaces},
arXiv:1103.5126 \cite{source}, ask in Section~8, ``Further remarks and open
problems,'' for a generalization to the Heckman--Opdam hypergeometric setting:
after recalling the extension of compact and noncompact spherical transforms
to hypergeometric functions associated with root systems, they call a
root-system version of Ramanujan's Master Theorem a natural question.  They
single out even multiplicity as a case where shift operators should reduce the
problem to a Euclidean $W$-invariant theorem.  This passage is on page~33 of
the PDF rebuilt from the cached arXiv source.

\section*{Supporting answer}

The same authors subsequently published
\emph{Ramanujan's Master theorem for the hypergeometric Fourier transform on
root systems}, arXiv:1211.0024 \cite{answer}.  Its abstract states that the
paper proves precisely such an analogue and that it generalizes the earlier
symmetric-space result to arbitrary positive multiplicity functions.

The exact result is Theorem~5.1, ``Ramanujan's Master Theorem for root
systems,'' on PDF pages 18--19.  It starts with a root system $\Sigma$ and a
positive multiplicity function $m$, forms the alternating normalized
Heckman--Opdam Jacobi series from a Hardy-class coefficient function, proves
its holomorphic convergence, represents its continuation by the
hypergeometric Fourier inversion integral, and identifies its hypergeometric
Fourier transform.  Section~6 proves the theorem.  Both the introduction and
the proof explicitly describe it as the extension of the earlier
Riemannian-symmetric-space theorem.

Thus arXiv:1211.0024 gives an author-known direct answer, and it is stronger
than the even-multiplicity subcase suggested in arXiv:1103.5126.

\newpage
\section*{Scope and evidence note}

This identification answers only the root-system extension.  The source's
separate broad remarks about higher-rank Bertram kernels and Paley--Wiener
characterization remain outside this status claim.

The official supporting PDF download was unavailable because the environment's
external-usage quota was exhausted.  The official arXiv abstract and primary
full-PDF text were inspected online, including Theorem~5.1 and Section~6; exact
links and locations are recorded in \texttt{supporting\_evidence.md}.  The
local \texttt{source\_paper.pdf} was compiled from the cached official arXiv
TeX source.

\begin{thebibliography}{9}
\bibitem{source}
G. Ólafsson and A. Pasquale,
\emph{Ramanujan's Master Theorem for Riemannian symmetric spaces},
J. Funct. Anal. 262 (2012), 2151--2180; arXiv:1103.5126.
\url{https://arxiv.org/abs/1103.5126}.

\bibitem{answer}
G. Ólafsson and A. Pasquale,
\emph{Ramanujan's Master theorem for the hypergeometric Fourier transform on
root systems}, arXiv:1211.0024 (2012).
\url{https://arxiv.org/abs/1211.0024}.
\end{thebibliography}

\end{document}
